%% ------------------------------------------------------------
%% Beyond the Scoresheet - one-page proposal (ACM two-column style)
%% Compile with pdfLaTeX in Overleaf (the default compiler).
%% ------------------------------------------------------------
\documentclass[sigconf,nonacm]{acmart}

% Remove the ACM copyright / conference boxes so it looks like the sample
\settopmatter{printacmref=false}
\setcopyright{none}
\renewcommand\footnotetextcopyrightpermission[1]{}
\pagestyle{plain}

\begin{document}

%% ---------------- TITLE & AUTHORS ----------------
\title{Modelling Football Match Events and Player Connections: Bayer Leverkusen 2023/24}
\subtitle{Project Codename: Beyond the Scoresheet}

\author{Meet Hemang Shah}
\affiliation{\institution{University of Victoria}}
\email{meetshah2@uvic.ca}

\author{Sakshi Pinkeshkumar Patel}
\affiliation{\institution{University of Victoria}}
\email{sakshipinkeshkumarp@uvic.ca}

\author{Ayush Bharatbhai Patel}
\affiliation{\institution{University of Victoria}}
\email{ayushbharatbhaipatel@uvic.ca}

%% ---------------- ABSTRACT ----------------
\begin{abstract}
This proposal tries to solve a data-modeling problem using a real, public sports dataset: showing how matches, teams, players, lineups, and timestamped match events are linked across all 34 Bundesliga games that Bayer 04 Leverkusen played during its historic, undefeated 2023–24 title-winning season.\cite{reuters2024}. Using StatsBomb's Open Data (competition ID~9, season
ID~281) \cite{statsbomb-open,statsbomb-list}, we study how relational and
semantic (RDF/OWL) representations preserve and expose the same
relationships --- such as a goal's link to its scorer, lineup context, and an
explicitly source-coded assist pass --- for a shared set of
relationship-oriented queries. The project's contribution is a way to show and rate things, not a way to predict or explain why the club won. 
\end{abstract}

\keywords{Sports Event Data Model, Relationship Modelling, Semantic Interoperability}

\maketitle

%% ---------------- 1 INTRODUCTION ----------------
\section{Introduction}
A football game involves many interacting factors, such as who participated,
where they played, and what they did, as well as shots, goals, cards,
replacements, and how everything worked together to decide the result. This
study examines the 2023--2024 Bayer Leverkusen Bundesliga season. Key domain
entities include matches, teams, players, positions, formations, and discrete
match events \cite{schemaorg}. The relationships between entities make this
domain difficult to represent: a goal cannot be separated from its scorer,
assist, passing sequence, pitch formation, or substitutions
\cite{statsbomb-match}. This project addresses the data-modelling problem of
how to represent the events of a football match --- the players and teams
involved, their actions, and the relationships that connect events across a
match or a season --- in a way that is precise, relationally complete, and
reusable across information needs \cite{cdf2025}, so that humans and
computational systems can reliably reconstruct who was involved in any given
event, trace how one event led to another \cite{provo}, and compare events.

%% ---------------- 2 MOTIVATION ----------------
\section{Motivation}
An integrated representation of football event data can link player, match,
lineup, and event information that is currently scattered among separate
records. This provides a unified view of player participation, positional
duties, substitutions, and subsequent actions across games, and enables more
complex queries that trace player--event interactions while keeping their time
and match context. Football analysts and scholars may find this useful for
quicker retrieval, comparison, and reuse of detailed match information.
Explicitly expressed relations also provide a cleaner base for automated
querying and knowledge extraction by computational systems. Hence, this study
explores ways to represent related football information while retaining key
contextual linkages.

\subsection{Motivating Example}
Let's say a Leverkusen player comes on as a replacement and creates a chance to score. To understand what that player did, several records must be linked: the
substitution, the player's position in the game, the next pass or shot, and
the game in which it happened. Current event records capture these facts in
isolation, but the relationships between them are not explicitly modelled as
one connected journey, which makes it difficult to track a player's path from
joining the match to creating an opportunity. This project encodes these
linkages so that such connected occurrences can be queried jointly, while
keeping their temporal and match-specific context.

%% ---------------- 3 TEAM JUSTIFICATION ----------------
\section{Team Justification}
Our team is composed of three people with complementary knowledge. All
participants have shared coursework experience in database architecture and
data administration. One member has football domain expertise, another has
hands-on experience with semantic web technologies, and another has hands-on
experience with applied artificial intelligence. This combination supports
domain analysis, data modelling, semantic representation, execution, and
evaluation of the project. Every member will be involved in all project phases,
and progress will be documented via our shared GitHub repository and commit
history.

%% ---------------- REFERENCES ----------------
\begin{thebibliography}{9}

\bibitem{reuters2024}
Reuters.
\newblock Leverkusen become first Bundesliga team to go all season without
  defeat.
\newblock 18 May 2024.

\bibitem{statsbomb-open}
StatsBomb.
\newblock StatsBomb Open Data, GitHub repository.
\newblock \url{https://github.com/statsbomb/open-data}

\bibitem{statsbomb-list}
StatsBomb.
\newblock Bundesliga 2023/24 match list, competition~9 / season~281.

\bibitem{schemaorg}
Schema.org.
\newblock SportsEvent.
\newblock \url{https://schema.org/SportsEvent}

\bibitem{statsbomb-match}
StatsBomb.
\newblock Bayer Leverkusen v Werder Bremen, 14 April 2024, event file, match
  3895302.

\bibitem{cdf2025}
Anzer et al.
\newblock Common Data Format (CDF) for Soccer Match Data.
\newblock arXiv:2505.15820, 2025.

\bibitem{provo}
W3C.
\newblock PROV-O: The PROV Ontology.
\newblock \url{https://www.w3.org/TR/prov-o/}

\end{thebibliography}

\end{document}
